\makeabstract
    {
    A Pneumatically Actuated Optical Mount for Enabling Beam Interchange
    }
    {
    Rafael G\'omez\textsuperscript{1}, David S. Hall\textsuperscript{1}
    }
    {
    \textsuperscript{1} Department of Physics \& Astronomy, Amherst College, Amherst, MA
    }
    {
    We imprint non-trivial topologies on trapped Rubidium-87 Bose-Einstein condensate samples using selected atomic transitions driven by laser beams. To refine our transition-driving ability, we are incorporating an improved projection beam; however, this setup relies on multiple distinct beams propagating to the atoms, in cycles 30 seconds apart, through different optics along the same stretch of beam path. To make beam interchange possible, this project aimed to develop a computer-controlled, translating optical mount.
    }
    {
    To move the optic, we preferred a pneumatic cylinder to alternatives like piezo actuators---which could introduce excess acoustic noise---and commercial flipper mounts, which do not fit within apparatus geometry. Hence, motion is propelled by an SMC 2-inch stroke cylinder and confined by four custom-machined parts in concert with a commercial guide rail. Air flow to the cylinder is controlled by a 3/2 solenoid valve, which takes in a high or low signal from a combined digital and analog soldered circuit. This circuit allows for manual or computer control. Through this system, the optical mount translates a mirror into apparatus beam path, directing one magneto-optical trap (MOT) beam into the Rb-87 cell before moving it out of beam path so that projection and imaging beams can proceed unobstructed. 
    }
    {
    Relief cuts and kinematic stops ensure the mount{\textquoteright}s movement is stable. The soldered circuit gives computer control and, after adaptation, controls coolant pumps on the optical table. Further, the custom parts have several tapped holes which interface with the optics, providing adaptability as we prepare the mount for experiments.
    }
    {
    This well-constrained, computer-controlled optical mount helps us meet the challenge posed by beam interchange. The mirror attached to the cylinder system directs beams to the atoms as intended, and computer control allows the mount{\textquoteright}s timing to align with the cycles of lasers in our experiment. Currently, the cylinder needs to be designed into the space near the Rb-87 cell. Together with our ongoing work to make beam paths more compact, this will help us switch to using a Spatial Light Modulator (SLM) to drive atomic transitions. Patterned light from the SLM will further enable us to imprint different excitations on the condensates and study their interactions.
    }

    \newpage
    